% !TeX encoding = UTF-8
\documentclass[variant=guiaresuelta,cover=institutional,mode=screen]{ua-engineering-v6-article}
% !TeX encoding = UTF-8
\UASetUniversity{Universidad Austral}
\UASetFaculty{Facultad de Ingeniería}
\UASetDepartment{Departamento de Computación}
\UASetCampus{Pilar}
\UASetCareer{Ingeniería Informática}
\UASetCommission{Comisión A}
\UASetProfessor{Equipo docente}
\UASetSemester{Segundo cuatrimestre 2026}
\UASetCourse{Sistemas Distribuidos}
\UASetDate{2026}


\UASetClassNumber{10}
\UASetClassTitle{Observabilidad distribuida}
\UASetDocKind{Guía resuelta}

\title{Clase 10 --- Guía resuelta: observabilidad, tracing, métricas y debugging}
\author{\UAProfessor}
\date{\UADate}

\begin{document}
\maketitle

\section{Criterio de lectura}
Las resoluciones son modelos de razonamiento, no recetas universales. Una buena respuesta separa hechos observados, hipótesis, evidencia confirmatoria o refutatoria y decisión. En producción, una explicación convincente debe poder ser contradicha por nueva evidencia.

\section{Resoluciones completas}

\begin{uaexercise}
Un equipo afirma que su sistema es observable porque dispone de un dashboard con CPU, memoria y cantidad de requests. Explique por qué eso constituye monitoreo útil, pero no demuestra observabilidad. Proponga una pregunta nueva sobre el comportamiento interno que esas métricas no podrían responder por sí solas.
\end{uaexercise}
\begin{uaanswer}
CPU, memoria y tráfico permiten detectar saturación o cambios agregados, pero no reconstruyen por qué una operación concreta falló ni cómo se propagó la latencia entre servicios. El monitoreo responde preguntas previstas de antemano; la observabilidad permite formular preguntas nuevas utilizando señales suficientemente ricas.

Una pregunta que ese dashboard no responde sería: ``¿qué dependencia y qué intento de retry explican que este usuario haya recibido dos cobros mientras la latencia global permanecía normal?''. Para responderla harían falta, como mínimo, un identificador de operación, logs estructurados correlacionados, la traza de la solicitud y evidencia del estado persistido.
\end{uaanswer}

\begin{uaexercise}
Para un checkout distribuido que invoca Auth, Inventory, Payments y Orders, diseñe un esquema mínimo de log estructurado. Distinga campos obligatorios para correlación, campos de negocio y campos que deberían evitarse o protegerse por privacidad.
\end{uaexercise}
\begin{uaanswer}
Campos de correlación: \texttt{timestamp}, \texttt{service.name}, \texttt{environment}, \texttt{trace\_id}, \texttt{span\_id}, \texttt{request\_id}, \texttt{operation}, \texttt{severity}, \texttt{status} y \texttt{duration\_ms}. Campos de negocio útiles: \texttt{order\_id}, \texttt{payment\_id}, \texttt{inventory\_reservation\_id}, tipo de operación e identificador seudonimizado del tenant.

No deberían registrarse contraseñas, tokens, números completos de tarjeta, secretos ni datos personales innecesarios. Los identificadores de usuario deben minimizarse o seudonimizarse y la retención debe respetar el propósito declarado. El log debe permitir unir hechos sin transformarse en una segunda base de datos sensible.
\end{uaanswer}

\begin{uaexercise}
Compare logs, métricas y trazas usando el mismo incidente: el 2\% de los checkouts comienza a tardar más de cuatro segundos. Indique qué pregunta responde mejor cada señal, qué no puede responder y cómo se combinan.
\end{uaexercise}
\begin{uaanswer}
Las métricas muestran alcance y tendencia: cuándo comenzó, qué porcentaje está afectado, qué región o versión concentra el problema y cómo cambiaron p95/p99, errores y saturación. No explican por sí solas una solicitud individual.

Las trazas muestran el camino de una muestra de checkouts y permiten localizar qué spans concentran el tiempo, qué dependencias se invocan y dónde aparecen retries. Pueden omitir casos por sampling.

Los logs aportan detalle discreto: códigos de error, decisiones, identificadores y cambios de estado. Sin correlación pueden convertirse en ruido. El flujo correcto es detectar con métricas, aislar patrones con trazas y confirmar detalles con logs, volviendo luego a las métricas para medir impacto.
\end{uaanswer}

\begin{uaexercise}
Una plataforma centraliza todos sus logs en un único pipeline. Durante una falla del pipeline, los servicios siguen funcionando, pero la organización pierde temporalmente visibilidad operacional. Analice por qué la propia plataforma de observabilidad debe tratarse como un sistema distribuido crítico. Proponga dos mecanismos de resiliencia.
\end{uaexercise}
\begin{uaanswer}
La telemetría tiene productores, buffers, red, procesadores, almacenamiento e índices; puede perder, duplicar, retrasar o reordenar datos igual que cualquier pipeline distribuido. Si falla, el servicio principal puede seguir disponible pero queda operativamente ciego, lo que eleva tiempo de detección y recuperación.

Dos mecanismos apropiados son: (1) buffering local durable con límites, backpressure y reenvío posterior; (2) rutas redundantes o almacenamiento temporal alternativo que permita detectar divergencia entre eventos emitidos y recibidos. También convienen SLI propios: porcentaje de eventos aceptados, lag de ingestión, tasa de descarte y completitud por origen.
\end{uaanswer}

\begin{uaexercise}
Diseñe los cuatro Golden Signals para un servicio de pagos. Para cada señal, especifique una métrica concreta, su unidad, una segmentación útil y una interpretación operacional.
\end{uaexercise}
\begin{uaanswer}
\textbf{Latencia}: p50/p95/p99 de autorización en milisegundos, segmentada por proveedor y región; una cola larga puede anticipar timeouts. \textbf{Tráfico}: autorizaciones por segundo, segmentadas por operación y tenant; permite contextualizar saturación. \textbf{Errores}: proporción de operaciones fallidas por causa y código; deben separarse rechazos de negocio de fallas técnicas. \textbf{Saturación}: utilización del pool de conexiones, profundidad de cola y concurrencia en vuelo; indica cuán cerca está el sistema de perder estabilidad.

Las cuatro señales se interpretan juntas: una tasa de error alta con tráfico bajo sugiere una falla distinta de una p99 alta acompañada por pool al 100\%.
\end{uaanswer}

\begin{uaexercise}
Un servicio reporta latencia media de 120 ms, p50 de 80 ms, p95 de 350 ms y p99 de 4.8 s. Explique qué historia cuenta esta distribución y por qué optimizar solo la media puede empeorar la experiencia de una minoría significativa.
\end{uaexercise}
\begin{uaanswer}
La mayoría de las solicitudes es rápida, pero existe una cola muy larga: al menos 1\% tarda cerca de segundos y 5\% supera ampliamente la experiencia típica. Esa minoría puede corresponder a una región, una dependencia, una clave caliente o un camino con retry.

Optimizar la media podría acelerar aún más el camino rápido sin tocar la cola. Para un sistema interactivo, el p99 de 4.8 s puede cruzar el timeout del cliente y generar retries, duplicación o abandono. Deben segmentarse percentiles y examinar trazas de la cola, no solo reducir el promedio.
\end{uaanswer}

\begin{uaexercise}
Aplique el método RED a un servicio HTTP y el método USE a un pool de conexiones de base de datos. Explique por qué ambos marcos son complementarios y no sustitutos.
\end{uaexercise}
\begin{uaanswer}
Para el servicio HTTP, RED observa \emph{Rate} (requests/s), \emph{Errors} (proporción por tipo) y \emph{Duration} (percentiles). Describe la experiencia y el comportamiento del servicio.

Para el pool, USE observa \emph{Utilization} (conexiones ocupadas/capacidad), \emph{Saturation} (espera o cola) y \emph{Errors} (timeouts, conexiones rotas). Describe el recurso que sostiene al servicio. RED puede mostrar p99 alta; USE ayuda a explicar si el pool está saturado. Uno mira demanda y resultado; el otro, capacidad del recurso.
\end{uaanswer}

\begin{uaexercise}
Un equipo agrega \texttt{user\_id}, \texttt{session\_id}, \texttt{order\_id} y la URL completa como etiquetas de cada métrica. Explique el problema de cardinalidad que aparece, sus costos y cómo rediseñaría la instrumentación sin perder capacidad diagnóstica.
\end{uaexercise}
\begin{uaanswer}
Cada combinación de etiquetas genera una serie temporal. Identificadores casi únicos multiplican series, memoria, índices, costo de almacenamiento y tiempo de consulta hasta volver inestable la plataforma de métricas.

La solución es reservar métricas para dimensiones acotadas: servicio, operación, región, status y versión. Los identificadores únicos deben ir en logs o trazas, donde pueden consultarse bajo demanda. La URL debe normalizarse por ruta lógica, por ejemplo \texttt{/orders/:id}. Exemplars pueden enlazar una observación métrica con una traza representativa sin convertir cada usuario en una serie.
\end{uaanswer}

\begin{uaexercise}
Diseñe una traza distribuida para el flujo \texttt{POST /checkout}. Incluya un span raíz y spans para Auth, Inventory, Payments, Orders y persistencia. Para cada span indique nombre, atributos relevantes, estado y relación padre-hijo.
\end{uaexercise}
\begin{uaanswer}
El span raíz puede llamarse \texttt{HTTP POST /checkout} y contener región, versión, tenant y status, sin datos sensibles. Sus hijos directos: \texttt{Auth.ValidateToken}, \texttt{Inventory.Reserve}, \texttt{Payments.Authorize} y \texttt{Orders.Create}. Cada servicio puede tener hijos internos, por ejemplo \texttt{DB INSERT order}.

Los atributos deben usar claves estables: \texttt{rpc.system}, \texttt{db.system}, \texttt{messaging.destination}, resultado y número de intento. El estado del span indica éxito o error técnico; los rechazos de negocio deben representarse sin confundirlos con falla de infraestructura. La jerarquía permite identificar camino crítico y paralelismo.
\end{uaanswer}

\begin{uaexercise}
El checkout publica un evento \texttt{OrderCreated} en Kafka y un consumidor asincrónico inicia el envío. Explique cómo propagar el contexto de traza a través del broker sin confundir causalidad asincrónica con una llamada RPC bloqueante.
\end{uaexercise}
\begin{uaanswer}
El productor inyecta el contexto de traza en headers del mensaje junto con \texttt{event\_id} y \texttt{correlation\_id}. El span de publicación representa la producción y el consumidor crea un nuevo span enlazado al contexto recibido. Cuando hay fan-out o retraso largo, un \emph{span link} expresa causalidad sin fingir una pila síncrona continua.

El offset no reemplaza al identificador causal: solo ubica el evento en una partición. El consumidor debe registrar topic, partición, offset, consumer group y edad del evento para diagnosticar lag, duplicación y replay.
\end{uaanswer}

\begin{uaexercise}
Compare head sampling, tail sampling y sampling adaptativo. Elija una política para un sistema con 100.000 requests por segundo donde interesa conservar todas las trazas con error y una muestra representativa de las exitosas.
\end{uaexercise}
\begin{uaanswer}
Head sampling decide al inicio y es barato, pero puede descartar una traza que luego termina en error. Tail sampling decide tras observar el resultado y puede conservar errores y latencias extremas, a costa de buffering y complejidad. El adaptativo modifica tasas según servicio, tráfico o anomalías.

Para este caso usaría una tasa pequeña de head sampling para cobertura uniforme más tail sampling para conservar 100\% de errores y trazas por encima de un umbral de latencia, con límites de capacidad. Además, revisaría sesgo por región y operación para evitar que caminos de bajo tráfico desaparezcan.
\end{uaanswer}

\begin{uaexercise}
Una traza muestra que Payments tarda 2.5 s. Los logs de Payments no contienen errores y sus métricas de CPU son normales. Formule al menos tres hipótesis alternativas y describa qué evidencia adicional buscaría para distinguirlas.
\end{uaexercise}
\begin{uaanswer}
Hipótesis 1: dependencia externa lenta; buscar subspans de red, DNS, TLS y proveedor. Hipótesis 2: espera por pool de conexiones; revisar saturación, tiempo de adquisición y cola. Hipótesis 3: bloqueo o contención en base de datos; examinar locks, queries y wait events. También puede haber pausa de runtime, rate limiting o timeout/retry interno.

La ausencia de error y CPU normal no prueba salud. La investigación debe descomponer los 2.5 s en tiempos de espera y ejecución, comparar trazas rápidas/lentas y buscar una dimensión común.
\end{uaanswer}

\begin{uaexercise}
Construya una cadena causal completa donde una base de datos lenta provoque: acumulación de requests, timeouts, retries, agotamiento del pool de conexiones y errores 500 en el borde. Distinga causa inicial, amplificadores, síntomas intermedios y síntoma visible al usuario.
\end{uaexercise}
\begin{uaanswer}
Causa inicial: una operación de base aumenta su latencia por lock, plan regresivo o I/O. Síntoma intermedio: las conexiones permanecen ocupadas más tiempo y crece la espera. Amplificador 1: timeouts del servicio generan retries. Amplificador 2: varias capas reintentan, multiplicando tráfico. Amplificador 3: el pool se agota y operaciones sanas también esperan.

Finalmente, la API supera su deadline y devuelve 500/504. El usuario ve checkout fallido; la primera alerta puede aparecer en el borde, aunque la causa original esté en una query. La evidencia debe unir p99, wait time del pool, número de intentos, spans de DB y cambios recientes.
\end{uaanswer}

\begin{uaexercise}
Durante un incidente, una persona afirma: ``la causa raíz es Payments porque ahí aparece el primer span rojo''. Explique por qué esa conclusión puede ser prematura. Proponga un procedimiento de hipótesis y refutación para investigar sin caer en correlación superficial.
\end{uaexercise}
\begin{uaanswer}
El span rojo indica dónde se observó un error, no necesariamente dónde se originó. Payments puede fallar porque recibió un deadline agotado, porque Inventory entregó datos inválidos o porque una dependencia externa no respondió.

Procedimiento: fijar el síntoma y ventana; comparar trazas afectadas y sanas; formular hipótesis alternativas; identificar predicciones observables de cada una; buscar subspans, métricas y logs; intentar refutar primero; revisar cambios y segmentaciones; reproducir o inyectar la condición de manera controlada. La explicación final debe cubrir mecanismo causal y no solo coincidencia temporal.
\end{uaanswer}

\begin{uaexercise}
Una saga de compra contiene eventos \texttt{StockReserved}, \texttt{PaymentAuthorized}, \texttt{ShipmentRequested} y compensaciones. Diseñe la observabilidad necesaria para reconstruir una instancia completa de saga, aun cuando los eventos lleguen tarde, duplicados o fuera de orden.
\end{uaexercise}
\begin{uaanswer}
Cada saga necesita \texttt{saga\_id}; cada evento, \texttt{event\_id}, tipo, versión, productor, tiempo de evento, tiempo de ingestión y causation ID. Los servicios registran transición anterior/nueva, intento e idempotency key. Las trazas usan links entre producción y consumo.

Una vista de auditoría ordena por causalidad y versión, no solo por timestamp. Debe mostrar duplicados descartados, compensaciones, lag y estado terminal. Así puede distinguirse ``evento tardío'' de ``paso nunca ejecutado'' y reconstruir qué supo cada servicio.
\end{uaanswer}

\begin{uaexercise}
Analice el siguiente incidente inspirado en un caso real: un cambio de configuración en el sistema de logs provoca pérdida de una parte de la telemetría durante varias horas, mientras el servicio principal continúa atendiendo tráfico. ¿Qué SLI propios debería tener la plataforma de observabilidad? ¿Cómo detectaría pérdida silenciosa y cómo diseñaría retención local o buffering?
\end{uaexercise}
\begin{uaanswer}
SLI: proporción de eventos recibidos respecto de emitidos por origen; lag de ingestión; tasa de descarte; cobertura de hosts; disponibilidad de consulta; completitud de campos. La pérdida silenciosa se detecta con contadores monotónicos o secuencias por productor, heartbeats de telemetría y reconciliación entre bytes/eventos emitidos y aceptados.

Los agentes mantienen spool local durable con límites, prioridad por severidad y backpressure. Al recuperar el backend reenvían preservando identificadores para deduplicar. Deben definirse políticas de descarte explícitas: preferir perder debug antes que auditoría, y alertar cuando el buffer se aproxima al límite.
\end{uaanswer}

\begin{uaexercise}
Diseñe una alerta para un SLO de 99.9\% de requests exitosas en 30 días. Compare una alerta por umbral instantáneo de error con una alerta basada en burn rate del error budget. Explique cuál reduce mejor ruido sin ocultar incidentes importantes.
\end{uaexercise}
\begin{uaanswer}
Un umbral como ``error rate >1\% durante cinco minutos'' ignora tamaño del budget, duración y ventana; puede alertar por picos irrelevantes o no alertar por una degradación lenta pero costosa.

El burn rate mide cuántas veces más rápido se consume el budget respecto de lo sostenible. Una estrategia multiventana combina una señal rápida de burn muy alto con otra lenta de burn moderado. Reduce ruido porque vincula alerta con riesgo real de incumplir el SLO y permite diferenciar incidente urgente de tendencia que requiere intervención planificada.
\end{uaanswer}

\begin{uaexercise}
La organización no puede conservar todos los logs y trazas indefinidamente. Diseñe una política de retención por niveles que considere: datos recientes de alta resolución, agregados históricos, trazas con error, privacidad y costo. Declare qué capacidad diagnóstica se pierde.
\end{uaexercise}
\begin{uaanswer}
Ejemplo: 7 días de logs y trazas detalladas; 30 días de errores, seguridad y auditoría; 90 días de métricas agregadas; histórico más largo solo para SLI y capacidad. Las trazas exitosas se muestrean, mientras errores y p99 se retienen más. Datos sensibles se minimizan, cifran y eliminan según propósito.

Se pierde capacidad de reconstruir requests exitosas antiguas con detalle y de investigar incidentes detectados tarde. El trade-off debe documentarse y compensarse con agregados, snapshots de incidentes y exportación legal cuando corresponda.
\end{uaanswer}

\begin{uaexercise}
Se observa el siguiente patrón: p99 de API aumenta, CPU permanece estable, longitud de cola de Orders crece y el lag del consumer group también aumenta. Diseñe un plan de investigación de diez pasos que combine métricas, trazas, logs, cambios recientes y experimentos controlados.
\end{uaexercise}
\begin{uaanswer}
1) fijar ventana y regiones; 2) verificar impacto en SLI; 3) segmentar por endpoint/versión; 4) revisar tasa de llegada y capacidad de consumo; 5) comparar lag por partición; 6) examinar saturación de pools, I/O y dependencias; 7) buscar trazas lentas representativas; 8) correlacionar logs por trace/event ID; 9) revisar despliegues, configuración y skew de claves; 10) validar la hipótesis con una reducción controlada de carga o réplica adicional.

CPU estable descarta solo una clase de saturación. La cola y el lag sugieren que la capacidad efectiva de un consumidor o dependencia cayó. La intervención debe medir recuperación y evitar ``arreglar'' el síntoma sin confirmar el mecanismo.
\end{uaanswer}

\begin{uaexercise}
Aplique el marco $D=(P,F,G,S,T)$ al diseño de observabilidad para una plataforma de reservas global. Incluya fallas de la propia telemetría, garantía de correlación, estrategia de sampling, SLI/SLO, costo aceptado y un experimento de fault injection que valide la propuesta.
\end{uaexercise}
\begin{uaanswer}
$P$: explicar y operar reservas que atraviesan regiones, pagos y proveedores. $F$: pérdida/retraso de telemetría, relojes distintos, sampling, cardinalidad, duplicación y caída del backend. $G$: correlación extremo a extremo de operaciones críticas, detección de degradación y SLI confiables aun con fallas parciales.

$S$: contexto W3C propagado por RPC y eventos; logs estructurados; RED/USE; tail sampling de errores/p99; spooling local; SLO de éxito y latencia; burn-rate alerts. $T$: costo de instrumentación, almacenamiento, privacidad y overhead. Experimento: cortar el collector de una región e introducir latencia en Payments; verificar que no se pierdan eventos críticos, que el buffer recupere y que la cadena causal siga reconstruible.
\end{uaanswer}

\end{document}
